\chapter{Software de adquisición SDR del sensor}

\section{Propósito del capítulo}

El propósito de este capítulo es describir, de forma técnica y verificable, la
arquitectura del código que corre directamente en el sensor de monitoreo espectral
remoto. El sensor está compuesto por una Raspberry Pi 5, una HackRF One y una
tarjeta auxiliar que integra comunicación LTE, receptor GPS y un multiplexor para
selección de antena.

La sección del proyecto descrita en este capítulo no corresponde a toda la
plataforma. El sistema completo también incluye componentes externos, como
backend, frontend, bases de datos y servicios de visualización. Este capítulo se
concentra en el código embebido del sensor: la parte que recibe instrucciones de
la plataforma, controla el hardware de radiofrecuencia, adquiere muestras,
procesa señales y devuelve resultados.

El propósito general del código es convertir al sensor físico en un nodo autónomo
de medición espectral. Para ello, el software debe recibir órdenes del servidor,
configurar la HackRF One, seleccionar la antena adecuada, capturar señales IQ,
procesarlas localmente y enviar a la plataforma central productos útiles como PSD,
métricas de modulación, audio demodulado, estado del dispositivo y coordenadas
GPS.

De manera específica, el software del sensor busca cumplir los siguientes
objetivos:

\begin{itemize}
  \item Controlar la adquisición de radio mediante la HackRF One, ajustando
  frecuencia central, tasa de muestreo, ganancias, amplificador y puerto de
  antena.

  \item Ejecutar mediciones espectrales en modo \textit{realtime} y en campañas
  programadas desde el servidor.

  \item Procesar muestras IQ crudas para estimar densidad espectral de potencia
  mediante métodos como Welch y banco de filtros polifase.

  \item Demodular señales AM o FM cuando el servidor solicita audio en tiempo
  real.

  \item Aplicar etapas de limpieza y corrección para mejorar la calidad de las
  mediciones, incluyendo balance IQ y remoción del pico DC.

  \item Mantener una calibración de frecuencia que compense desviaciones del
  oscilador de la HackRF.

  \item Coordinar la comunicación con el backend, el envío de resultados, los
  reintentos ante fallos de red y el reporte de estado del sensor.

  \item Gestionar recursos físicos del sensor, como LTE, GPS, GPIO y multiplexor
  de antena.
\end{itemize}

El capítulo describe el software actualmente implementado en el repositorio. La
capa de control se encuentra principalmente en \texttt{orchestrator.py},
\texttt{campaign\_runner.py}, \texttt{status.py}, \texttt{retry\_queue.py},
\texttt{functions.py}, \texttt{cfg.py} y \texttt{utils/}. La capa de adquisición y
procesamiento RF se encuentra en \texttt{rf/rf.c} y \texttt{rf/libs/*.c}. Los
servicios de posicionamiento y conectividad se encuentran en \texttt{gps-lte/},
mientras que el audio en tiempo real se articula con \texttt{server\_webrtc.py}.
Por tanto, las funcionalidades descritas no son una especificación abstracta, sino
una lectura organizada del código del sensor.


\section{Contexto dentro del sistema general}

El software del sensor hace parte del sistema de monitoreo espectral multisensor
remoto. Dentro de la arquitectura general, el sensor actúa como un nodo de
adquisición y procesamiento local, mientras que la plataforma central se encarga
de coordinar tareas, almacenar resultados, visualizar información y permitir la
interacción con el usuario.

La arquitectura del sensor se puede entender como una cadena de trabajo: el
servidor decide qué se debe medir, el orquestador traduce esa instrucción a una
configuración de radio, el motor RF adquiere y procesa las muestras, y finalmente
los resultados regresan al servidor.

La división principal del sistema es la siguiente:

\begin{itemize}
  \item \textbf{Orquestación en Python}: gestiona la relación con el backend,
  consulta instrucciones, programa campañas, controla estados del sistema, envía
  comandos al motor RF y prepara los resultados para publicarlos.

  \item \textbf{Motor RF en C}: controla la HackRF One, recibe muestras IQ,
  administra buffers de alta velocidad, ejecuta procesamiento DSP, calcula PSD,
  demodula audio y devuelve resultados al orquestador.

  \item \textbf{Servicios auxiliares}: atienden funciones de soporte como GPS,
  LTE, estado del sistema, reintentos de datos no enviados, WebRTC para audio y
  control del multiplexor de antena.
\end{itemize}

La decisión de separar Python y C responde a una necesidad práctica. Python es
adecuado para coordinación, comunicación HTTP, manejo de estados y programación
de tareas. C es adecuado para la parte sensible a latencia: captura de muestras,
buffers, FFT, filtrado, PSD y demodulación.

El sensor opera con varios servicios o modos complementarios:

\begin{itemize}
  \item \textbf{PSD en tiempo real}: el backend solicita una configuración y el
  sensor responde con una estimación espectral de la banda capturada.

  \item \textbf{PSD y audio en tiempo real}: además de la PSD, el sensor demodula
  FM o AM y genera tramas de audio reproducibles por la plataforma.

  \item \textbf{Campañas}: el backend programa mediciones periódicas; el sensor
  guarda la configuración, agenda ejecuciones y sube resultados en cada
  intervalo.

  \item \textbf{Calibración en frecuencia}: el sensor estima el error de sintonía
  y aplica una corrección para que las mediciones queden mejor alineadas en
  frecuencia.

  \item \textbf{GPS y LTE}: el sensor mantiene conectividad celular y reporta
  ubicación geográfica al servidor.

  \item \textbf{Estado del dispositivo}: se reportan variables de salud como
  conectividad, almacenamiento, temperatura, memoria y actividad del software.
\end{itemize}

\subsection{Arquitectura funcional del software SDR}

La ruta principal de una medición puede resumirse como una cadena de control y una
cadena de datos. La cadena de control inicia en el backend, pasa por el
orquestador Python y llega al motor RF mediante mensajería local. La cadena de
datos inicia en la HackRF One, produce muestras IQ, ejecuta procesamiento digital
de señales en C y retorna productos resumidos hacia la plataforma.

% Requiere en el preambulo:
% \usepackage{tikz}
% \usetikzlibrary{arrows.meta,positioning}
\begin{figure}[htbp]
\centering
\begin{tikzpicture}[
  scale=0.92,
  transform shape,
  font=\tiny,
  external/.style={
    draw,
    rounded corners=2pt,
    very thick,
    align=center,
    text width=2.35cm,
    minimum height=0.88cm,
    fill=blue!8
  },
  orchbox/.style={
    draw,
    rounded corners=4pt,
    very thick,
    align=center,
    minimum width=3.35cm,
    minimum height=4.25cm,
    fill=blue!4
  },
  motorbox/.style={
    draw,
    rounded corners=4pt,
    very thick,
    align=center,
    minimum width=3.35cm,
    minimum height=4.75cm,
    fill=green!10
  },
  submain/.style={
    draw,
    rounded corners=2pt,
    thick,
    align=center,
    text width=2.78cm,
    minimum height=0.48cm,
    fill=blue!10
  },
  subrf/.style={
    draw,
    rounded corners=2pt,
    thick,
    align=center,
    text width=2.78cm,
    minimum height=0.48cm,
    fill=green!16
  },
  aux/.style={
    draw,
    rounded corners=2pt,
    thick,
    align=center,
    text width=2.55cm,
    minimum height=0.72cm,
    fill=orange!13
  },
  flow/.style={-Latex, very thick},
  ctrl/.style={-Latex, very thick},
  hw/.style={Latex-Latex, very thick},
  side/.style={-Latex, thick},
  parallel/.style={-Latex, thick, dashed},
  edgeLabel/.style={midway, fill=white, inner sep=1pt, align=center}
]

\node[external] (backend) at (-0.35,0) {\textbf{Backend / API}\\solicitudes, campañas\\resultados};

\node[orchbox] (orch) at (3.10,0) {};
\node[font=\tiny\bfseries, anchor=north] at (orch.north) {Orquestador Python};
\node[submain] (api) at (3.10,1.23) {\textbf{HTTP / API REST}\\consulta y publicación};
\node[submain] (services) at (3.10,0.45) {\textbf{Control de servicios}\\realtime, campañas, audio};
\node[submain] (cron) at (3.10,-0.33) {\textbf{Control de tiempos}\\cron de campañas};
\node[submain] (post) at (3.10,-1.11) {\textbf{Postproceso}\\remoción de pico DC};

\node[motorbox] (motor) at (7.25,0) {};
\node[font=\tiny\bfseries, anchor=north] at (motor.north) {Motor RF en C};
\node[subrf] (acq) at (7.25,1.42) {\textbf{Control de adquisición}\\configuración HackRF};
\node[subrf] (buffers) at (7.25,0.62) {\textbf{Lógica de buffers}\\recepción y drenaje IQ};
\node[subrf] (preiq) at (7.25,-0.18) {\textbf{Preproceso IQ}\\balance y corrección DC};
\node[subrf] (psd) at (7.25,-0.98) {\textbf{Procesamiento}\\PSD, audio y métricas};
\node[subrf] (omp) at (7.25,-1.78) {\textbf{Multihilo}\\OpenMP y tareas paralelas};

\node[external, fill=green!10] (hackrf) at (10.95,-0.35) {\textbf{HackRF One}\\sintonía, ganancias\\flujo IQ};
\node[aux] (status) at (-0.35,2.35) {\textbf{Servicio status}\\temperatura, CPU, RAM\\disco y ping};
\node[aux] (webrtc) at (3.10,2.75) {\textbf{Servicio WebRTC}\\audio paralelo\\RTP / ICE};
\node[aux] (gpio) at (10.95,1.35) {\textbf{GPIO / MUX}\\selección física\\de antena};

\draw[ctrl] (backend.east) -- (orch.west);
\draw[ctrl] (orch.east) -- (motor.west);
\draw[hw] (motor.east) -- ++(0.45,0) |- (hackrf.west);
\draw[flow] (psd.west) -- (post.east);

\draw[parallel] (orch.north) -- (webrtc.south);
\draw[parallel] (status.south) -- (backend.north);
\draw[side] (acq.east) -- ++(0.35,0) |- (gpio.west);
\draw[side] (gpio.south) -- (hackrf.north);

\node[align=center, font=\tiny] at (5.45,-3.05) {
Azul: plano de control/publicación \quad
Verde: adquisición y DSP\\
Naranja: servicios paralelos o auxiliares
};

\end{tikzpicture}
\caption{Diagrama de bloques de la arquitectura funcional del software SDR del sensor.}
\label{fig:arquitectura-software-sdr}
\end{figure}

En el diagrama, el orquestador Python se muestra como el proceso de coordinación:
consulta y publica por API REST, activa servicios de tiempo real, campañas y
audio, agenda campañas mediante cron y aplica el postproceso final de PSD. El
motor RF en C concentra el camino crítico de adquisición y DSP: configura la
HackRF, administra buffers, preprocesa IQ, estima PSD, produce métricas y usa
multihilo para las etapas de cómputo intensivo. El servicio WebRTC corre en
paralelo al orquestador y se activa cuando hay demodulación; toma audio
demodulado del camino RF y lo entrega mediante WebRTC usando RTP e ICE. El
servicio de estado también corre de forma continua y reporta métricas de salud
del sistema hacia la plataforma.

\begin{center}
\begin{tabular}{p{0.24\textwidth}p{0.66\textwidth}}
\toprule
\textbf{Bloque} & \textbf{Responsabilidad técnica} \\
\midrule
Backend & Publica solicitudes de medición, campañas, parámetros RF y recibe resultados. \\
Orquestador Python & Consulta tareas por API REST, coordina servicios, agenda campañas y aplica postproceso de PSD. \\
ZeroMQ IPC & Transporta comandos JSON locales entre Python y C usando un flujo estricto REQ/REP. \\
Motor RF C & Configura HackRF, controla adquisición, buffers, preproceso IQ, PSD, métricas, audio y multihilo. \\
libhackrf/USB & Ejecuta la comunicación de bajo nivel con la HackRF One y entrega muestras IQ intercaladas. \\
WebRTC/Opus & Corre en paralelo, se activa con demodulación y entrega audio mediante WebRTC, RTP e ICE. \\
Servicio status & Reporta temperatura, carga de CPU, RAM usada, espacio en disco y ping hacia la plataforma. \\
ShmStore & Comparte estado en \texttt{/dev/shm/persistent.json}, como \texttt{ppm\_error}, GPS y campañas. \\
GPS/LTE/GPIO & Mantiene ubicación, conectividad celular y control físico del multiplexor de antenas. \\
\bottomrule
\end{tabular}
\end{center}

Esta separación define una frontera importante: el sensor ejecuta adquisición,
procesamiento básico y control local; la plataforma central se encarga de
persistencia histórica, visualización, comparación institucional, reportes y
analítica de mayor nivel. La tabla siguiente resume esa distribución.

\begin{center}
\begin{tabular}{p{0.28\textwidth}p{0.62\textwidth}}
\toprule
\textbf{Nivel} & \textbf{Funciones principales} \\
\midrule
Sensor/edge & Adquisición IQ, PSD local, filtrado, demodulación AM/FM, GPS, estado, reintentos y audio. \\
Backend/cloud & Almacenamiento histórico, campañas, gestión de sensores, APIs, consolidación y reglas institucionales. \\
Frontend & Visualización, selección de parámetros, marcadores visuales, navegación de resultados y operación del usuario. \\
\bottomrule
\end{tabular}
\end{center}


\section{Desarrollo técnico}

El código puede leerse como un conjunto de bloques funcionales. Cada bloque cumple
una responsabilidad dentro del camino completo de la señal: desde la instrucción
del backend hasta la publicación del resultado.

\subsection{Adquisición}

El bloque de adquisición es el encargado de hablar con la HackRF One y convertir
una instrucción de medición en muestras IQ reales.

Conceptualmente, este bloque realiza cinco acciones:

\begin{enumerate}
  \item Recibe una configuración de medición enviada por el orquestador. Esa
  configuración contiene parámetros como frecuencia central, ancho de muestreo,
  ganancias, método de PSD, puerto de antena, filtro solicitado y modo de
  demodulación.

  \item Aplica la configuración sobre la HackRF One. Esto incluye sintonía de
  frecuencia, tasa de muestreo, ganancias de recepción y corrección de frecuencia
  cuando existe una calibración disponible.

  \item Inicia la recepción de muestras IQ crudas desde el dispositivo SDR.

  \item Deposita esas muestras en buffers circulares para separar la velocidad de
  llegada del hardware de la velocidad de procesamiento del software.

  \item Entrega bloques de muestras suficientemente grandes a los algoritmos de
  procesamiento, evitando que cada operación pesada interrumpa directamente la
  captura.
\end{enumerate}

Las muestras IQ representan la señal recibida en banda base compleja. La componente
\textit{I} corresponde a la parte en fase y la componente \textit{Q} a la parte en
cuadratura. En conjunto permiten conservar información de amplitud y fase, que es
necesaria para estimar el espectro, filtrar, demodular FM o demodular AM.

En términos operativos, una adquisición completa sigue el siguiente procedimiento:

\begin{enumerate}
  \item Validar que la máquina de estados global permita la medición solicitada.
  El sistema evita ejecutar adquisiciones simultáneas entre \texttt{REALTIME},
  \texttt{CAMPAIGN} y \texttt{KALIBRATING}.

  \item Construir una configuración JSON con frecuencia central, tasa de muestreo,
  RBW, ventana, ganancias, antena, filtro, modo de demodulación y error de
  frecuencia \texttt{ppm\_error} cuando está disponible.

  \item Enviar la configuración al motor RF por ZeroMQ y esperar una respuesta
  antes de emitir otro comando, porque el contrato local es REQ/REP.

  \item Seleccionar el puerto de antena mediante GPIO cuando el hardware auxiliar
  lo requiere.

  \item Configurar la HackRF One con \texttt{libhackrf}: frecuencia central,
  \textit{sample rate}, LNA, VGA, amplificador y corrección de frecuencia.

  \item Iniciar el flujo USB de muestras IQ, llenar buffers circulares y vigilar
  condiciones de saturación, ausencia de dispositivo o error de streaming.

  \item Cerrar de forma controlada la captura, liberar recursos temporales y
  devolver al orquestador una respuesta JSON con resultados o error.
\end{enumerate}

\subsection{Parámetros SDR configurables}

Los parámetros de medición llegan desde el backend, son validados por el
orquestador y finalmente son interpretados por el parser del motor RF. La tabla
resume los parámetros más relevantes para adquisición y procesamiento.

\begin{longtable}{p{0.23\textwidth}p{0.18\textwidth}p{0.43\textwidth}}
\toprule
\textbf{Parámetro} & \textbf{Unidad/formato} & \textbf{Uso en el sensor} \\
\midrule
\endhead
Frecuencia central & Hz & Sintonía nominal de la HackRF. Puede corregirse internamente con \texttt{ppm\_error}. \\
Sample rate & muestras/s & Ancho de adquisición instantáneo y base para calcular resolución frecuencial. \\
Span & Hz & Banda observada alrededor de la frecuencia central; en la práctica queda limitado por el sample rate útil. \\
FFT size / \texttt{nperseg} & muestras & Número de puntos por segmento para Welch o tamaño de FFT en procesamiento espectral. \\
RBW & Hz & Resolución aproximada de frecuencia; se relaciona con \texttt{sample\_rate}/\texttt{FFT size}. \\
Overlap Welch & fracción & Porcentaje de solape entre segmentos para estabilizar la PSD. \\
Ventana & texto & Tipo de ventana espectral: Hamming, Hann, Blackman, Kaiser, Tukey, entre otras. \\
Método PSD & texto & Selección entre Welch y banco de filtros polifase cuando se requiere mayor selectividad. \\
LNA/VGA & dB configurados & Ganancias de recepción de la HackRF; afectan nivel relativo y riesgo de saturación. \\
Amplificador RF & booleano & Habilita o deshabilita el amplificador frontal de la HackRF. \\
Puerto de antena & entero & Selección física del camino RF mediante el multiplexor de antenas. \\
Filtro digital & Hz inicio/fin & Recorta la región de interés dentro de la banda realmente capturada. \\
Demodulación & AM/FM/ninguna & Activa el camino de audio en tiempo real. \\
Cooldown & segundos & Pausa solicitada entre capturas o ciclos de consulta. \\
ppm\_error & ppm & Corrección de frecuencia estimada durante calibración. \\
\bottomrule
\end{longtable}

\subsection{Span, sample rate, FFT y RBW}

En el sensor, \textit{sample rate}, span, tamaño de FFT y RBW no son parámetros
independientes. La tasa de muestreo define el ancho instantáneo máximo que puede
observarse en una sola captura. El tamaño de FFT o \texttt{nperseg} define cuántos
puntos se usan para dividir esa banda en contenedores de frecuencia. Como regla
práctica:

\[
\mathrm{RBW} \approx \frac{\mathrm{sample\_rate}}{\mathrm{FFT\ size}}
\]

Una RBW menor mejora la separación entre componentes cercanas, pero exige más
muestras por segmento, más memoria temporal, más cómputo FFT y normalmente mayor
latencia. Una RBW mayor reduce costo computacional y mejora la respuesta temporal,
pero disminuye la capacidad de distinguir señales próximas. Por esta razón, el
motor RF no debe reducir unilateralmente la tasa de muestreo o el tamaño de FFT
solicitados para ahorrar recursos: esos parámetros hacen parte de la medición
definida por el servidor.

\subsection{Uso de buffers}

La lógica de buffers es central en el sensor. La HackRF entrega datos de forma
continua y rápida, mientras que la PSD, el filtrado o la demodulación requieren
procesamiento por bloques. Por eso el motor RF usa buffers circulares:

\begin{itemize}
  \item Un buffer principal recibe muestras IQ para estimación espectral.
  \item Un buffer de audio conserva muestras necesarias cuando se activa
  demodulación AM o FM.
  \item El software drena estos buffers cuando ya hay suficientes muestras para
  construir una respuesta.
\end{itemize}

Esta estructura permite que el sensor atienda procesos en tiempo real, como PSD y
audio, sin tratar cada muestra de manera aislada. La idea pedagógica es similar a
una banda transportadora: el hardware deposita muestras, el procesamiento toma
paquetes completos y el orquestador recibe resultados ya elaborados.

\subsection{Preprocesamiento}

El preprocesamiento agrupa las operaciones que preparan la señal antes de extraer
productos finales. Su objetivo es reducir artefactos propios del hardware y
mejorar la calidad de la medición.

\subsubsection{Balance IQ}

En un receptor SDR real, las ramas I y Q no siempre quedan perfectamente
balanceadas. Puede haber diferencias de ganancia, pequeños desplazamientos DC o
correlación no deseada entre ambas componentes. Si estos efectos no se corrigen,
pueden aparecer artefactos en el espectro o degradación en la demodulación.

El balance IQ busca compensar esos errores antes del cálculo de PSD o del filtrado.
De forma conceptual, el algoritmo:

\begin{itemize}
  \item remueve desplazamientos promedio en las ramas I y Q;
  \item compara la energía relativa de ambas ramas;
  \item corrige diferencias de amplitud;
  \item reduce correlaciones espurias entre I y Q.
\end{itemize}

No se trata de cambiar la información de la señal recibida, sino de disminuir
imperfecciones introducidas por la cadena de recepción.

\subsubsection{Corrección de componente DC}

La componente DC aparece cerca del centro de la banda base y suele estar asociada
a efectos internos del receptor. En el flujo del sensor se atiende en dos niveles:

\begin{itemize}
  \item En el preproceso IQ se reduce el offset de las muestras antes de calcular
  PSD o demodular.
  \item En el postproceso de PSD se repara el pico central visible en el espectro
  cuando aparece como un artefacto de medición.
\end{itemize}

Esta separación es importante: una cosa es limpiar la señal cruda y otra es pulir
el producto espectral que se enviará al servidor.

\subsection{Procesamiento}

El bloque de proceso transforma muestras IQ en productos útiles para la plataforma.
En este repositorio se implementan tres familias principales de procesamiento:
estimación de PSD, demodulación de audio y filtrado de señales.

\subsubsection{Estimación de PSD}

La PSD, o densidad espectral de potencia, permite observar cómo se distribuye la
energía de la señal dentro de un rango de frecuencias. Es el producto principal
para monitoreo espectral, porque permite identificar ocupación de banda, niveles
relativos de potencia, portadoras y posibles emisiones.

El sensor contempla métodos como:

\begin{itemize}
  \item \textbf{Welch}: divide la señal en segmentos, aplica ventanas, calcula FFT
  y promedia resultados para obtener una estimación estable del espectro.

  \item \textbf{Banco de filtros polifase}: mejora la separación entre canales y
  ayuda a controlar fugas espectrales cuando se requiere una estimación más
  selectiva.
\end{itemize}

En la implementación, el método Welch es el camino principal para entregar una PSD
estable en tiempo real y en campañas. El banco de filtros polifase se reserva para
casos donde se requiere mejor separación espectral o menor fuga entre canales,
asumiendo un mayor costo computacional. La configuración incluye tamaño de
segmento, solape, tipo de ventana, tasa de muestreo y método PSD; con estos datos
se genera el eje de frecuencia y el vector de potencia relativo.

La PSD es una de las tareas más costosas del sistema, porque requiere operar sobre
grandes bloques de muestras y ejecutar transformadas de Fourier. Por eso el motor
RF está escrito en C y aprovecha una lógica de trabajo multihilo, planes FFTW y
buffers reutilizables por hilo. El uso de cachés locales reduce la sobrecarga de
planificación FFT, pero exige cuidado con memoria dinámica, OpenMP y almacenamiento
local de hilo.

\subsubsection{Escala de potencia y unidades}

El resultado logarítmico que produce el sensor debe interpretarse como una escala
relativa al rango digital del receptor, es decir, como dBFS o potencia relativa,
no como una medición absoluta en dBm. En este proyecto no se realizó una
calibración de potencia de la cadena RF completa. Por tanto, no se conoce con
certeza la tensión equivalente de saturación del ADC, las atenuaciones internas de
la HackRF, las pérdidas de cables, el factor de antena, la ganancia real de cada
etapa ni la respuesta en frecuencia del conjunto antena-cable-SDR.

Esta aclaración es crítica: aunque el sistema puede mostrar diferencias relativas
entre señales, detectar ocupación y comparar niveles dentro de una misma
configuración, no debe reportarse como instrumento calibrado en dBm, dB$\mu$V/m,
mV/m, W/m$^2$ o dBW/m$^2$/MHz sin una calibración externa. Para convertir la
escala actual a unidades absolutas se requeriría caracterizar el sistema con un
generador RF calibrado, niveles conocidos, varias frecuencias, varias ganancias,
registro de pérdidas y una curva de corrección con incertidumbre asociada.

\subsubsection{Demodulación FM y AM}

Cuando el backend solicita audio en tiempo real, el sensor no solo calcula PSD.
También toma las muestras IQ, extrae la información modulada y la convierte en
audio reproducible.

En FM, la información está principalmente en la variación de fase o frecuencia
instantánea. El proceso conceptual consiste en observar cómo cambia la fase entre
muestras consecutivas, convertir esa variación en una señal de audio y acondicionar
el resultado para reproducción.

En AM, la información está principalmente en la envolvente de la señal. El proceso
conceptual consiste en calcular la magnitud de la señal compleja, separar las
variaciones útiles de la portadora y normalizar el audio resultante.

Después de demodular, el audio se convierte a un formato de transmisión. El sensor
genera tramas de audio, las codifica con Opus y las entrega al servicio local que
publica el flujo mediante WebRTC. Esto permite que la plataforma escuche la señal
con baja latencia sin que el backend tenga que procesar IQ crudo.

\subsubsection{Filtrado}

El filtrado permite concentrarse en una parte de la banda capturada y rechazar
contenido fuera del rango de interés. El servidor puede solicitar un rango de
frecuencia y el motor RF aplica un filtro pasabanda antes de calcular la PSD o de
preparar la señal.

Desde una perspectiva conceptual, el filtro funciona como una ventana selectiva:
conserva la región que interesa y atenúa lo que queda por fuera. En el código se
apoya en procedimientos de dominio frecuencial y validaciones para que el rango
solicitado sea coherente con la banda realmente capturada por la HackRF.

Actualmente el filtro se expresa como un intervalo de frecuencias de inicio y fin
dentro de la banda capturada. Si el rango solicitado cae por fuera de los límites
de Nyquist definidos por la frecuencia central y la tasa de muestreo, el parser
ajusta el intervalo al rango físicamente observable. Filtros pasabajo o pasaaltos
pueden representarse como casos particulares de ese intervalo, pero una
caracterización formal de orden, ancho de transición y atenuación fuera de banda
debería documentarse como mejora futura si se requiere trazabilidad metrológica.

\subsubsection{Trabajo multihilo}

El sistema separa tareas para mantener capacidad de respuesta:

\begin{itemize}
  \item Un flujo atiende la llegada de muestras desde la HackRF.
  \item Otro flujo procesa bloques para PSD.
  \item Cuando hay audio, un hilo adicional drena muestras, demodula y envía
  tramas codificadas.
  \item El orquestador permanece disponible para comunicarse con el backend y
  manejar cambios de modo.
\end{itemize}

Esta organización evita que una tarea pesada, como una PSD de alta resolución,
bloquee completamente la captura o la transmisión de audio.

La lógica de multitarea y multihilo usada para paralelizar regiones de cómputo
dentro del motor RF se apoya en la librería OpenMP para C. Esta herramienta
permite distribuir operaciones repetitivas, como acumulaciones, compensación IQ y
segmentos de PSD, entre varios hilos de ejecución de forma relativamente sencilla,
sin tener que implementar manualmente toda la administración de hilos mediante
interfaces de menor nivel.

El uso de OpenMP aporta varias ventajas al desarrollo del sensor. En primer lugar,
facilita la paralelización de operaciones repetitivas o costosas, como cálculos
sobre bloques de muestras, procesamiento espectral o tareas que pueden dividirse
entre núcleos de CPU. En segundo lugar, reduce la complejidad del código frente a
una implementación completamente manual con hilos, bloqueos y sincronización
explícita. Finalmente, es una solución ampliamente soportada en entornos de
compilación en C, lo que permite mantener una implementación portable y fácil de
integrar dentro del motor RF. No obstante, OpenMP se usa con cautela: el código
evita depender de referencias inseguras a almacenamiento local de hilo dentro de
regiones paralelas y controla la cantidad de hilos para no saturar la Raspberry Pi
5 durante PSD de alta resolución o audio simultáneo.

\subsection{Postprocesamiento}

El postprocesamiento corresponde a la pulida final aplicada a los resultados antes
de enviarlos al servidor. En este sistema, la etapa más importante es la corrección
del pico DC en la PSD.

\subsubsection{Remoción del pico DC en PSD}

Aunque el preproceso reduce offset en las muestras IQ, en la PSD puede quedar un
pico artificial cerca del centro de la banda. Ese pico no siempre representa una
emisión real; con frecuencia es un artefacto propio del receptor SDR.

El algoritmo de postproceso analiza el comportamiento estadístico de la PSD como
una señal con variabilidad local. En lugar de borrar siempre un número fijo de
puntos, estima la región afectada y decide cómo reconstruirla.

Conceptualmente, el procedimiento realiza:

\begin{itemize}
  \item detección de la región central afectada usando simetría y cambios locales
  en la pendiente;
  \item análisis de ocupación espectral alrededor del centro;
  \item validaciones con descriptores estadísticos como media, mediana,
  dispersión y comportamiento de la cola alta;
  \item ampliación de la ventana de reparación cuando la zona tiene baja
  ocupación espectral;
  \item reconstrucción local mediante interpolación lineal o ajuste polinomial.
\end{itemize}

El resultado enviado al backend conserva una PSD más limpia. Cuando aplica, el
software también puede conservar información de diagnóstico sobre la corrección,
como la ventana reparada y el modo de reconstrucción usado.

\subsection{Calibración en frecuencia}

La calibración en frecuencia busca compensar la desviación entre la frecuencia
nominal solicitada y la frecuencia real a la que queda sintonizado el receptor.
En SDR de bajo costo, esta desviación puede aparecer por tolerancias del oscilador,
temperatura o condiciones de operación.

El sensor usa una referencia conocida para estimar el error de frecuencia. En el
caso de emisiones FM comerciales, una referencia útil es el tono piloto estéreo de
19 kHz que aparece después de demodular FM. Si una estación FM fuerte contiene ese
tono, el sensor puede usarlo como punto de comparación.

El flujo conceptual es:

\begin{enumerate}
  \item Capturar una porción amplia de la banda donde se esperan emisoras FM.
  \item Identificar candidatos con potencia suficiente.
  \item Bajar cada candidato a banda base y demodular FM.
  \item Buscar el tono piloto alrededor de 19 kHz en la señal demodulada.
  \item Comparar la posición observada con la posición esperada.
  \item Estimar un error de frecuencia y expresarlo como factor de corrección.
\end{enumerate}

Cuando hay varias emisiones candidatas, el sistema puede usar más de una referencia
para evitar depender de una sola estación. Las emisiones más confiables reciben
mayor peso, por ejemplo si tienen mejor potencia o una detección más clara del tono
piloto. Con esa información se obtiene una estimación más robusta del error.

El resultado de la calibración se guarda como un error de frecuencia, normalmente
expresado en partes por millón. Luego, cuando el backend solicita una frecuencia
central, el orquestador incluye ese factor en la configuración enviada al motor RF.

El motor usa la corrección para ajustar la sintonía de la HackRF. De esta manera,
la frecuencia física de adquisición queda más alineada con la frecuencia nominal
que la plataforma espera visualizar.

\subsection{Orquestador}

El orquestador es la lógica de coordinación del sensor. Su función no es procesar
IQ de forma intensiva, sino dirigir el trabajo de todos los bloques.

Sus responsabilidades principales son:

\begin{itemize}
  \item Consultar al backend para saber si hay solicitudes en tiempo real.
  \item Consultar y sincronizar campañas programadas.
  \item Validar configuraciones antes de enviarlas al motor RF.
  \item Enviar comandos al motor RF usando mensajería local ZeroMQ.
  \item Recibir resultados de PSD y métricas desde el motor RF.
  \item Aplicar postproceso sobre los resultados cuando corresponde.
  \item Formatear y publicar datos hacia el backend mediante HTTP REST.
  \item Controlar el inicio y parada del servicio de audio WebRTC.
  \item Mantener estados globales como \texttt{IDLE}, \texttt{REALTIME},
  \texttt{CAMPAIGN}, \texttt{KALIBRATING} y \texttt{ERROR}.
  \item Gestionar memoria compartida local mediante \texttt{ShmStore}; el archivo
  usado es \texttt{/dev/shm/}\\\texttt{persistent.json}. Allí se conservan
  parámetros como \texttt{ppm\_error}, coordenadas GPS o configuraciones de
  campaña.
\end{itemize}

Las campañas son mediciones programadas por el servidor. El orquestador consulta
las campañas activas, selecciona la que corresponde ejecutar, guarda sus parámetros
y agenda ejecuciones periódicas. Cada ejecución captura PSD y envía el resultado
al backend. Si la red falla, el resultado puede quedar en una cola local para ser
reenviado posteriormente.

En modo \textit{realtime}, el servidor entrega una configuración inmediata. El
orquestador la envía al motor RF, espera la respuesta, aplica postproceso y sube
el resultado. Mientras el backend siga enviando una solicitud válida, el sensor
mantiene este modo activo.

Cuando la solicitud incluye demodulación, el orquestador activa también el camino
de audio. En ese caso, el motor RF demodula AM o FM y genera tramas Opus; el
servicio WebRTC se encarga de hacer disponible ese audio para la plataforma.

Antes o durante ciertos ciclos de operación, el orquestador puede solicitar al
motor RF una calibración de frecuencia. El resultado se guarda localmente y se usa
en adquisiciones posteriores.


\section{Entradas, salidas e interfaces}

El software del sensor se comunica con diferentes elementos internos y externos.
Por esta razón, se puede describir en términos de entradas, salidas e interfaces.

\subsection{Entradas}

Las principales entradas del sistema son:

\begin{itemize}
  \item Configuraciones de medición enviadas desde el backend.
  \item Parámetros de campañas programadas.
  \item Solicitudes de PSD en tiempo real.
  \item Solicitudes de PSD con audio en tiempo real.
  \item Parámetros de frecuencia central, span, ganancias, filtro, puerto de
  antena y modo de demodulación.
  \item Muestras IQ crudas entregadas por la HackRF One.
  \item Datos GPS recibidos desde el módulo de posicionamiento.
  \item Estado de conectividad entregado por el servicio LTE.
  \item Parámetros persistentes almacenados localmente, como el error de
  frecuencia en partes por millón.
\end{itemize}

\subsection{Salidas}

Las principales salidas del sistema son:

\begin{itemize}
  \item Estimaciones de PSD enviadas al backend.
  \item Métricas de modulación y resultados derivados del procesamiento.
  \item Audio demodulado en tiempo real cuando se solicita AM o FM.
  \item Tramas de audio codificadas con Opus.
  \item Coordenadas GPS del sensor.
  \item Estado general del dispositivo.
  \item Información de conectividad, almacenamiento, temperatura, memoria y
  actividad del software.
  \item Resultados de calibración en frecuencia.
  \item Datos almacenados localmente para reintento cuando hay fallos de red.
\end{itemize}

Estas salidas tienen formatos y alcances diferentes. Para evitar ambigüedad, la
tabla siguiente resume qué producto genera el sensor y qué significado tiene.

\begin{longtable}{p{0.22\textwidth}p{0.24\textwidth}p{0.40\textwidth}}
\toprule
\textbf{Salida} & \textbf{Formato} & \textbf{Observación técnica} \\
\midrule
\endhead
PSD & JSON con eje de frecuencia y potencia relativa & Producto principal de monitoreo. La escala logarítmica es relativa/dBFS si no existe calibración de potencia. \\
Audio & Tramas Opus/WebRTC & Disponible solo cuando se solicita demodulación AM o FM en tiempo real. \\
GPS & Campos de estado y coordenadas & Debe asociarse a la medición con estado de fix cuando esté disponible. \\
Estado & JSON HTTP & Incluye variables de salud como conectividad, temperatura, memoria, almacenamiento y actividad. \\
Calibración & \texttt{ppm\_error} & Corrección de frecuencia persistida en memoria compartida. No equivale a calibración de potencia. \\
Reintentos & JSON local/cola & Resultados no enviados por falla de red, para posterior publicación. \\
Logs & Archivos en \texttt{Logs/} & Evidencia operativa para diagnóstico y soporte. \\
\bottomrule
\end{longtable}

\subsection{Interfaces y protocolos}

Los protocolos y procedimientos principales que conectan los bloques del sensor
son los siguientes:

\begin{longtable}{p{0.30\textwidth}p{0.62\textwidth}}
\toprule
\textbf{Elemento} & \textbf{Uso conceptual en el sensor} \\
\midrule
\endhead
HTTP REST & Comunicación entre sensor y backend para consultar instrucciones y publicar datos. \\
ZeroMQ & Canal local de comandos entre el orquestador Python y el motor RF en C. \\
libhackrf/USB & Interfaz de control y adquisición de muestras desde la HackRF One. \\
WebRTC & Entrega de audio demodulado hacia cliente/plataforma con baja latencia. \\
Opus & Codificación de audio antes de enviarlo al servicio WebRTC. \\
GPIO & Control de líneas físicas para multiplexor de antena y soporte de hardware. \\
PPP/LTE & Conectividad celular del sensor hacia la red. \\
GPS/NMEA & Obtención y conversión de coordenadas geográficas. \\
Welch & Método de estimación de PSD por segmentos y promediado. \\
Banco polifase & Método de estimación espectral con mejor control de separación entre canales. \\
OpenMP & Paralelización de tareas de cómputo dentro del motor RF desarrollado en C. \\
\bottomrule
\end{longtable}

\subsection{Contrato local Python-C}

La comunicación entre el orquestador y el motor RF usa ZeroMQ mediante el
transporte \texttt{ipc://} y el endpoint local \texttt{/tmp/rf\_engine}. Aunque la clase Python se llama
\texttt{ZmqPairController}, el socket usado por Python es \texttt{REQ} y el socket
del motor C es \texttt{REP}. Por tanto, el flujo es estrictamente secuencial: por
cada solicitud debe recibirse una respuesta antes de emitir el siguiente comando.
Esta restricción complementa la máquina de estados global y evita adquisiciones
concurrentes sobre la misma HackRF.

\begin{longtable}{p{0.20\textwidth}p{0.32\textwidth}p{0.34\textwidth}}
\toprule
\textbf{Comando lógico} & \textbf{Parámetros principales} & \textbf{Respuesta esperada} \\
\midrule
\endhead
PSD/realtime & Frecuencia central, sample rate, RBW, ventana, método PSD, ganancias, antena y filtro & JSON con frecuencia, PSD, estado y metadatos de adquisición. \\
PSD con audio & Parámetros PSD más \texttt{demodulation=am/fm} & PSD y activación del flujo Opus/WebRTC para audio. \\
Campaña & Configuración guardada por campaña, ventana temporal e intervalo & Resultado PSD por ejecución y publicación o reintento. \\
Calibración de frecuencia & \texttt{calibrate=true} y configuración RF base & JSON con estado de calibración completa y \texttt{ppm\_error}. \\
Filtro & \texttt{start\_freq\_hz}, \texttt{end\_freq\_hz} & Procesamiento limitado al rango válido dentro de la banda capturada. \\
Parada/reset & Cambio de modo o solicitud sin demodulación & Cierre del camino de audio o retorno a estado de espera según corresponda. \\
\bottomrule
\end{longtable}

Los errores relevantes de esta interfaz incluyen JSON inválido, timeout de
ZeroMQ, respuesta tardía, HackRF no disponible, frecuencia fuera del rango
soportado, tasa de muestreo inválida, saturación USB, falla de streaming, filtro
fuera de banda, error de demodulación o falla de publicación HTTP. Ante timeout,
el controlador Python recrea el socket para descartar respuestas tardías y
restaurar el estado REQ/REP.


\section{Decisiones de diseño o implementación}

Una de las decisiones principales de diseño fue dividir el software del sensor en
dos capas: una capa de orquestación en Python y un motor RF en C. Esta separación
permite aprovechar las ventajas de cada lenguaje. Python facilita la comunicación
con el backend, el manejo de estados, la programación de campañas y la integración
con servicios auxiliares. C permite ejecutar tareas de adquisición y procesamiento
digital de señales con mayor control sobre memoria, buffers y rendimiento.

Otra decisión importante fue usar buffers circulares para desacoplar la velocidad
de llegada de las muestras desde la HackRF de la velocidad de procesamiento del
software. Esta estrategia permite atender flujos de datos continuos y evitar que
operaciones costosas, como la estimación de PSD o la demodulación, interrumpan la
captura de muestras.

También se decidió usar ZeroMQ como canal local de comunicación entre el
orquestador Python y el motor RF en C. Esto permite separar responsabilidades,
mantener independencia entre procesos y enviar comandos o recibir resultados sin
integrar todo el sistema en un único ejecutable. La elección de REQ/REP mantiene
un contrato simple y controlado: una solicitud de adquisición, calibración o audio
debe cerrarse con una respuesta antes de iniciar otra operación RF.

Para la entrega de audio en tiempo real, se optó por demodular localmente la señal
en el sensor, codificar el audio con Opus y publicarlo mediante WebRTC. Esta
decisión evita enviar IQ crudo al backend para tareas de escucha y reduce la carga
de procesamiento de la plataforma central.

En cuanto a procesamiento paralelo, se decidió usar OpenMP para implementar la
lógica de multitarea y multihilo dentro del motor RF. Esta elección se tomó porque
OpenMP facilita la paralelización de regiones de código en C, tiene soporte amplio
en compiladores comunes y permite aprovechar mejor los núcleos disponibles de la
Raspberry Pi 5 sin aumentar excesivamente la complejidad del código.

Finalmente, se incluyeron etapas de preprocesamiento y postprocesamiento para
mejorar la calidad de los productos enviados al backend. Entre estas etapas se
encuentran el balance IQ, la reducción de componente DC en las muestras y la
remoción del pico DC en la PSD.


\section{Problemas presentados y ajustes realizados}

Durante el diseño del software fue necesario considerar varios problemas propios
de un sistema SDR remoto.

Uno de los primeros problemas es la presencia de imperfecciones en la cadena de
recepción. En un receptor SDR real, las ramas I y Q pueden presentar diferencias
de ganancia, desplazamientos DC o correlación no deseada. Para mitigar estos
efectos se incorporó una etapa de balance IQ y reducción de offset antes del
procesamiento espectral o la demodulación.

Otro problema identificado es la aparición de un pico artificial cerca del centro
de la PSD. Este pico no necesariamente corresponde a una emisión real, sino que
puede ser un artefacto de medición asociado al receptor. Para atenderlo se incluyó
una etapa de postprocesamiento que analiza la región central de la PSD, estima la
zona afectada y reconstruye localmente el resultado mediante interpolación lineal
o ajuste polinomial.

También se contemplaron fallos de red asociados a la conectividad LTE. Por esta
razón, el orquestador debe manejar reintentos, estados internos y almacenamiento
temporal de resultados cuando el envío al backend no puede realizarse de forma
inmediata.

La desviación de frecuencia de la HackRF también fue considerada como un problema
relevante. Debido a tolerancias del oscilador, temperatura o condiciones de
operación, la frecuencia real de adquisición puede diferir de la frecuencia nominal
solicitada. Para corregir este efecto se implementó una lógica de calibración en
frecuencia basada en referencias presentes en emisiones FM comerciales.

Además, la operación en tiempo real exige que el sistema no bloquee la adquisición
mientras procesa datos. Para esto se ajustó la arquitectura mediante buffers
circulares, separación de tareas y paralelización con OpenMP dentro del motor RF.

En la operación real también deben considerarse errores de adquisición y de
periféricos: HackRF no detectada, saturación del bus USB, overflow de buffers,
frecuencia o sample rate inválidos, antena no disponible, falla de GPIO, pérdida
de LTE, ausencia de fix GPS, timeout de ZeroMQ, caída del backend o temperatura
alta de la Raspberry Pi. El software atiende parte de estas condiciones con
estados, reintentos, logs, recreación de sockets y retorno a \texttt{IDLE}; otras
condiciones requieren diagnóstico operativo o criterios de aceptación definidos
por el despliegue.


\section{Validación, pruebas o evidencias}

La validación del software del sensor puede entenderse desde el flujo completo de
funcionamiento. De forma resumida, el ciclo operativo del sistema es el siguiente:

\begin{enumerate}
  \item El sensor se identifica ante la plataforma mediante su dirección MAC y su
  estado local.

  \item El orquestador consulta si hay una tarea \textit{realtime} o una campaña
  activa.

  \item Si hay tarea, se construye una configuración de radio y se envía al motor
  RF.

  \item El motor RF configura la HackRF, selecciona antena si aplica y captura IQ.

  \item Las muestras pasan por preprocesamiento para mejorar su condición de
  medición.

  \item El bloque de proceso calcula PSD, aplica filtros y, si se solicita,
  demodula audio.

  \item El postproceso corrige artefactos finales, especialmente el pico DC en la
  PSD.

  \item El orquestador empaqueta el resultado y lo envía al backend.

  \item Si el envío falla, el dato puede quedar almacenado localmente para
  reintento.

  \item Los servicios auxiliares mantienen GPS, LTE, estado del sensor y audio
  según el modo activo.
\end{enumerate}

Este flujo permite validar que las principales etapas del sensor están conectadas:
recepción de instrucciones, adquisición IQ, procesamiento, postprocesamiento,
publicación de resultados y soporte de servicios auxiliares.

Como evidencias funcionales del sistema se consideran:

\begin{itemize}
  \item La capacidad de recibir configuraciones desde el backend.
  \item La configuración efectiva de la HackRF One.
  \item La adquisición de muestras IQ.
  \item La generación de PSD mediante métodos como Welch o banco de filtros
  polifase.
  \item La activación de demodulación AM o FM cuando se solicita audio.
  \item La entrega de audio codificado mediante Opus y publicado por WebRTC.
  \item El reporte de estado del sensor, GPS, conectividad y variables de salud.
  \item La aplicación de corrección de frecuencia cuando existe una calibración
  disponible.
  \item El manejo de reintentos ante fallos de comunicación con el backend.
\end{itemize}

Para que la validación sea repetible, se propone registrar cada prueba con entrada,
procedimiento, salida esperada y criterio de aceptación. La siguiente matriz
resume pruebas mínimas para el software del sensor.

\begin{longtable}{p{0.22\textwidth}p{0.34\textwidth}p{0.30\textwidth}}
\toprule
\textbf{Prueba} & \textbf{Procedimiento} & \textbf{Criterio de aceptación} \\
\midrule
\endhead
Conexión HackRF & Iniciar motor RF y enviar configuración PSD básica. & El motor responde JSON válido y no reporta dispositivo ausente. \\
PSD con señal conocida & Inyectar o sintonizar una portadora conocida. & El pico aparece en la frecuencia esperada dentro de la tolerancia de calibración. \\
RBW/FFT & Ejecutar capturas con distinto RBW y mismo sample rate. & La separación de bins cambia según sample rate/FFT sin degradar parámetros solicitados. \\
Demodulación FM/AM & Solicitar audio con señal compatible. & Se activa WebRTC/Opus y se recibe audio sin bloquear la PSD. \\
GPS/LTE & Ejecutar reporte de estado en campo. & Se reporta conectividad y coordenadas cuando existe fix válido. \\
Campaña & Programar ventana de campaña y esperar ejecución. & Se genera medición en el intervalo esperado y se publica o encola. \\
Caída de red & Interrumpir LTE/backend durante envío. & El resultado queda en cola local y se reintenta posteriormente. \\
Calibración de frecuencia & Ejecutar calibración con emisora FM fuerte. & Se obtiene \texttt{ppm\_error} razonable o se rechaza la calibración si no hay referencia. \\
Carga térmica & Ejecutar PSD de alta resolución por tiempo prolongado. & El sensor mantiene operación y reporta temperatura/estado para diagnóstico. \\
\bottomrule
\end{longtable}

\section{Trazabilidad frente a requisitos funcionales}

La tabla siguiente relaciona funciones esperadas en un sistema de monitoreo
espectral con el estado del software descrito en este capítulo. La columna de
estado diferencia lo implementado en el sensor, lo parcialmente cubierto y lo que
debe tratarse como recomendación o desarrollo futuro.

\begin{longtable}{p{0.27\textwidth}p{0.18\textwidth}p{0.39\textwidth}}
\toprule
\textbf{Funcionalidad} & \textbf{Estado} & \textbf{Comentario técnico} \\
\midrule
\endhead
Frecuencia central & Implementado & Configurable desde backend y aplicada por HackRF con corrección \texttt{ppm\_error}. \\
Sample rate/span hasta 20 MHz & Implementado & La captura instantánea depende de la tasa de muestreo útil de la HackRF. \\
RBW y FFT & Implementado & RBW derivada de sample rate y tamaño FFT/\texttt{nperseg}. \\
Ventanas y Welch & Implementado & Soporta ventanas y solape para estabilizar PSD. \\
Banco polifase & Implementado / parcial & Disponible como método PSD alternativo con mayor costo de CPU. \\
VBW/suavizado & Parcial & Puede aproximarse por promediado; falta formalizar como parámetro metrológico separado. \\
Filtros digitales & Parcial & Implementado como rango pasabanda validado; falta caracterización formal LP/HP/PB. \\
Demodulación AM/FM & Implementado & Audio local con codificación Opus y publicación WebRTC. \\
GPS & Implementado & Se reportan coordenadas; se recomienda documentar fix, precisión y HDOP si están disponibles. \\
LTE/reintentos & Implementado & Envío HTTP y cola local ante fallos de red. \\
Unidades dBm/dB$\mu$V/m & Pendiente & Requieren calibración de potencia, antena, pérdidas y curva de corrección. \\
Escala dBFS relativa & Implementado & Es la interpretación válida actual de los niveles logarítmicos del sensor. \\
Marcadores, OBW y potencia de canal & Parcial/futuro & Pueden calcularse sobre PSD; debe definirse si se ejecutan en sensor, backend o frontend. \\
Sweep mayor a 20 MHz & Futuro & Requiere barrido por ventanas, solape, tiempo de permanencia y unión espectral. \\
Almacenamiento IQ & Parcial/futuro & Debe formalizar formato \texttt{.cs8}/metadatos si se requiere evidencia forense. \\
NTP/timestamping & Parcial & Se usan tiempos del sistema; se recomienda documentar fuente, zona horaria y tolerancia. \\
Monitoreo del sensor & Implementado / parcial & Estado, temperatura, memoria, almacenamiento y conectividad reportados para diagnóstico. \\
Seguridad HTTPS/TLS & Implementado / parcial & La URL base usa HTTPS; falta documentar autenticación, certificados y manejo de credenciales. \\
\bottomrule
\end{longtable}


\section{Aprendizajes y transferencia de conocimiento}

El desarrollo del software del sensor deja varios aprendizajes importantes para
la continuidad del proyecto.

En primer lugar, se evidencia la conveniencia de separar las tareas de
orquestación y las tareas de procesamiento intensivo. Python resulta adecuado para
la integración con servicios, consultas HTTP, manejo de estados y coordinación
general. C resulta más adecuado para las tareas críticas de adquisición,
procesamiento DSP, manejo de buffers y demodulación.

En segundo lugar, el uso de buffers circulares es fundamental en sistemas de
adquisición continua. En un sensor SDR, el hardware puede entregar muestras de
forma sostenida y a alta velocidad, por lo que el software debe estar preparado
para recibir datos sin bloquearse mientras ejecuta procesamiento por bloques.

En tercer lugar, se confirma la importancia del preprocesamiento y el
postprocesamiento. La señal entregada por un SDR real puede contener artefactos
como desbalance IQ, offset DC o picos centrales en la PSD. Por tanto, el software
no solo debe adquirir y calcular espectros, sino también aplicar correcciones que
mejoren la calidad de los productos entregados al backend.

Otro aprendizaje importante es que las funcionalidades de tiempo real dependen
tanto del procesamiento local como de la comunicación con la plataforma. Aunque el
motor RF procese los datos rápidamente, la experiencia de visualización también
depende de la latencia de red, la publicación al backend y la actualización de la
interfaz.

Finalmente, el uso de OpenMP muestra que es posible mejorar el aprovechamiento de
los recursos de cómputo disponibles sin aumentar excesivamente la complejidad del
código. Esta herramienta permite introducir paralelización en C de forma más
sencilla que una implementación manual completa basada en hilos.


\section{Limitaciones}

El sistema presenta varias limitaciones que deben tenerse en cuenta para su
operación y evolución futura.

Una primera limitación está asociada a la latencia de comunicación con el servidor.
Aunque el sensor pueda adquirir y procesar señales en tiempo real, la tasa de
actualización observada en la plataforma siempre estará limitada por los tiempos de
envío, recepción, procesamiento en backend y visualización en frontend. Por tanto,
el tiempo real del sistema no depende únicamente de la velocidad del motor RF, sino
también de la infraestructura de comunicación y de la plataforma central.

Otra limitación corresponde al span configurable. El sistema trabaja con un límite
de span de hasta 20 MHz, asociado al ancho de banda práctico de adquisición de la
HackRF One y a la configuración implementada en el software. Actualmente no se
incluye una funcionalidad de \textit{sweep} que permita recorrer secuencialmente
varias porciones del espectro para observar un ancho de banda mayor al que permite
una sola captura directa del hardware.

También existe una limitación en la calibración en frecuencia. La estrategia
implementada depende de detectar emisiones FM comerciales y, en particular,
referencias como el tono piloto estéreo de 19 kHz. Por esta razón, en zonas
aisladas o en lugares donde no existan emisoras FM suficientemente fuertes o
claras, el sensor puede no contar con referencias adecuadas para ejecutar la
calibración. Alternativas futuras incluyen una referencia RF externa, GPSDO,
beacons conocidos, estaciones patrón, comparación contra analizador certificado o
calibración periódica en laboratorio.

Otra limitación importante corresponde a la calibración en potencia. Las HackRF
One usadas en el sistema no son instrumentos calibrados en potencia absoluta y en
este proyecto no se realizó una calibración de amplitud de la cadena RF. Por esta
razón, al calcular el logaritmo de la potencia de la PSD, el resultado no debe
interpretarse como dBm ni como dB absolutos referenciados. La interpretación
correcta actual es relativa al rango digital de la adquisición, es decir, dBFS o
nivel relativo. No se conoce con certeza la tensión de saturación equivalente del
ADC, los rangos internos de tensión, las atenuaciones de la HackRF, las pérdidas
del cableado, el factor de antena ni la ganancia real del conjunto. Para obtener
mediciones robustas de potencia absoluta, sería necesario realizar una calibración
externa con señales de referencia, instrumentación adecuada y una curva de
corrección que relacione los niveles medidos por el sensor con niveles reales de
potencia.

Finalmente, al tratarse de un sistema remoto basado en conectividad LTE, el envío
de resultados puede verse afectado por cobertura, disponibilidad de red, pérdida
de paquetes o variaciones de latencia. Aunque el software contempla reintentos y
almacenamiento local, la operación continua depende de la calidad de la conexión
en el sitio de despliegue.


\section{Recomendaciones específicas}

A partir del análisis del software del sensor, se proponen las siguientes
recomendaciones específicas:

\begin{itemize}
  \item Mantener la separación entre orquestación en Python y procesamiento RF en
  C, ya que esta división permite conservar claridad arquitectónica y rendimiento
  en las tareas críticas.

  \item Documentar con mayor detalle las interfaces entre el orquestador y el
  motor RF, especialmente los mensajes ZeroMQ, los parámetros de configuración y
  los formatos de respuesta.

  \item Conservar el uso de buffers circulares y revisar periódicamente sus
  tamaños, tiempos de drenaje y condiciones de saturación para evitar pérdidas de
  muestras.

  \item Continuar usando OpenMP para paralelizar tareas de procesamiento en C,
  evaluando cuidadosamente qué secciones del código realmente se benefician de la
  paralelización.

  \item Implementar en el futuro una funcionalidad de \textit{sweep} si se requiere
  observar un span mayor a 20 MHz. Esta funcionalidad debería coordinar capturas
  sucesivas en diferentes frecuencias centrales y reconstruir una vista espectral
  más amplia.

  \item Desarrollar una metodología de calibración en potencia si el sistema se va
  a usar para mediciones absolutas. Esta calibración debería usar un generador RF
  calibrado, niveles conocidos, varias frecuencias, varias combinaciones de
  ganancia, registro de pérdidas de cable, factor de antena, incertidumbre y una
  curva de corrección para convertir dBFS a unidades referenciadas como dBm.

  \item Complementar la calibración en frecuencia con otros métodos de referencia
  para escenarios donde no existan emisoras FM comerciales disponibles o donde el
  tono piloto de 19 kHz no pueda detectarse de manera confiable, por ejemplo
  GPSDO, señal patrón externa, beacons conocidos o comparación con un instrumento
  certificado.

  \item Medir y caracterizar la latencia extremo a extremo del sistema, desde la
  adquisición en el sensor hasta la visualización en la plataforma, para definir
  con precisión el límite real de actualización en tiempo real.

  \item Mantener registros de diagnóstico sobre corrección de pico DC, calibración
  de frecuencia, fallos de red y estado del dispositivo, con el fin de facilitar
  soporte técnico y depuración.

  \item Formalizar timestamping y georreferenciación por medición, incluyendo
  fuente de tiempo, sincronización NTP, zona horaria, formato ISO 8601, estado de
  fix GPS, número de satélites y precisión cuando estén disponibles.

  \item Definir responsabilidades para marcadores, OBW, potencia de canal,
  umbrales y exportación: qué se calcula en el sensor, qué se deriva en backend y
  qué se visualiza o ajusta en frontend.

  \item Evaluar periódicamente el desempeño térmico y de recursos de la Raspberry
  Pi 5, especialmente durante tareas de PSD de alta resolución, demodulación de
  audio y operación multihilo.
\end{itemize}
